Составить список требований, предъявляемых к разрабатываемому веб-сайту (в соответствии с вариантом).
Требования должны делиться на следующие категории:
\begin{itemize}
\item Функциональные.
\begin{itemize}
\item Требования пользователей сайта.
\item Требования владельцев сайта.
\end{itemize}
\item Нефункциональные.
\end{itemize}
Требования необходимо оформить в соответствии с шаблонами RUP
(документ SRS --- Software Requirements Specification).
Для каждого из требований нужно указать его атрибуты
(в соответствии с методологией RUP),
а также оценить и аргументировать приблизительное количество часов,
требующихся на реализацию этого требования.

Для функциональных требований нужно составить UML UseCase-диаграммы, описывающие реализующие их прецеденты использования.

Отчёт по лабораторной работе должен содержать:
\begin{itemize}
\item Документ Software Requirements Specification, содержащий список требований к сайту.
\item UseCase-диаграммы прецедентов использования, реализующих функциональные требования.
\item Выводы по работе.
\end{itemize}

Описание веб-сайта:
\begin{quote}
  vk.com: vk is the largest european social network with more than 100 million active users. our goal is to keep old friends, ex-classmates, neighbors and colleagues in touch - \url{https://vk.com/}
\end{quote}
